\documentclass[12pt]{amsart} 
\usepackage{stackengine,scalerel}

\input{./latex-preamble/cm-preamble.tex}

%\graphicspath{ {./diagrams/} }

\newtheorem{lemma}{Lemma}

\begin{document}

{Hartshorne} \hfill {2-14-2025}

\bigskip\bigskip

\begin{center}

{\sc Chapter II.3, Week 4}

\end{center}

\begin{center}

  {Noah Parker} %\footnote{Collaborated with Kalev Martinson and Frank :).}}
    
\end{center}

\bigskip

% \begin{enumerate}
% 	\item[3.9.] {
%       \emph{The Topological Space of a Product.}
%       Recall that in the category of varieties, the Zariski topology on the product of two varieties is not equal to the product topology.
%       Now, we see that in the category of schemes, the underlying point set of a product of schemes is not even the product set.
%       \begin{enumerate}
%         \item Let $k$ be a field, and let $\A^1_k=\Spec k[x]$ be the affine line over $k$.
%           Show that $\A_k^1\times_{\Spec k}\A_k^1\cong \A_k^2$, and show that the underlying point set of the product is not the product of the underlying point set of the factors (even if $k$ is algebraically closed).
%         \item Let $k$ be a field with $s$ and $t$ indeterminates over $k$.
%           Then $\Spec k(s)$, $\Spec k(t)$, and $\Spec k$ are all one point spaces.
%           Describe the product scheme $\Spec k(s)\times_{\Spec k}\Spec k(t)$.
%       \end{enumerate}
%     }
% \end{enumerate}
% 
% \begin{proof}
%   (a) We have that $k[x]\otimes_k k[x]\cong k[x,y]$ as $k$-algebras, so $\A_k^1\times_{\Spec k}\A_k^1\cong \A_k^2$ as schemes over $k$.
%   Fix $k=\CC$, and consider the generic points of $\A_k^2$.
%   %If there were a bijection between $\A_k^2$ and the set theoretic product, then this would give us a bijection of generic points, and thus a bijection of irreducible components.
%   %There is an injection of point sets $f:\A_k^1\times\A_k^1\to \A_k^2$ given by $(\p,\q)\mapsto \p^e + \q^e$.
%   A bijection of sets would give us a bijection of generic points, and thus a bijection of irreducible polynomials.
%   This is clearly false, so that $\A_k^2\neq \A_k^1\times \A_k^1$ as sets.
% 
%   (b) Fix $A=k(s)\otimes_k k(t)$, so that the product scheme is given by $X=\Spec A$.
%   Consider the ideal $\p=(s\otimes 1 +1\otimes t)\subset A$.
%   Clearly $\p\neq (0)$.
%   If $\p=(1)$, then we have
%   \begin{align*}
%     1\otimes 1 &= (s\otimes 1 + 1\otimes t)\left(\sum_{i=1}^n \frac{f_i(s)}{p_i(s)}\otimes \frac{g_i(t)}{q_i(t)}\right) \\
%       &= \sum_{i=1}^n\frac{sf_i(s)}{p_i(s)}\otimes \frac{g_i(t)}{q_i(t)} +
%       \frac{f_i(s)}{p_i(s)}\otimes \frac{tg_i(t)}{q_i(t)}.
%   \end{align*}
% \end{proof}

\begin{enumerate}
	\item[3.10.] {
      \emph{Fibres of a morphism.}
      \begin{enumerate}
        \item If $f:X\to Y$ is a morphism and $y\in Y$ a point, show that $\ssp(X_y)$ is homeomorphic to $f^{-1}(y)$ with the induced topology.
        \item Let $X=\Spec k[s,t]/(s-t^2)$, $Y=\Spec k[s]$, and $f:X\to Y$ the morphism defined by $s\mapsto s$.
          If $y\in Y$ is a point $a\in K$ with $a\neq 0$, show that the fibre $X_y$ consists of two points, with residue field $k$.
          If $y\in Y$ corresponds to $0\in k$, show that the fibre $X_y$ is a nonreduced one-point scheme.
          If $\eta$ is the generic point of $Y$, show that $X_\eta$ is a one-point scheme, whose residue field is an extension of degree two of the residue field of $\eta$. (Assume $k$ algebraically closed.)
      \end{enumerate}
    }
\end{enumerate}

\begin{proof}
  (a) We begin by considering the affine case, $Y=\Spec A$, $X=\Spec B$, and $f^\sharp(Y)=\vphi:A\to B$.
  An identical method to that of Ex. 2.1 gives us the following lemma.
  \begin{lemma}
    Let $A$ be a ring and $\phi:A\to S^{-1}A$ the inclusion into its localization.
    Then $\phi^*:\Spec (S^{-1}A)\to \Spec A$ is a homeomorphism onto its image in $\Spec A = X$.
    We denote this image $S^{-1} X$.
  \end{lemma}
  % % Additionally, $\vphi_y:A_y\to (B_A)_y$ gives us the following diagram.
  % % \[
  % %   \begin{tikzcd}
  % %     A\ar[r,"\vphi"]\ar[d,"\phi_A"'] & B_A\ar[d,"\phi_B"] \\
  % %     A_y\ar[r,"\vphi_y"'] & (B_A)_y
  % %   \end{tikzcd}
  % % \] 
  % % Taking $\Spec$ gives us
  % % Since $\vphi_y$ is a local ring, $\vphi_y$ factors to a map
  % % \[
  % %   \ol{\vphi_y}:B_A
  % % \] 
  The next lemma is easy, but worth proving.
  \begin{lemma}
    Let $f:A\to B$ be a ring morphism, $X:= \Spec A$, $Y:=\Spec B$, and $S$ a multiplicative subset of $A$.
    Then denoting $S^{-1}B =f(S)^{-1}B\cong B_A\otimes_A S^{-1}A$ and identifying spectra as in the prior lemma, 
    \[
      S^{-1}f^*:\Spec (S^{-1}B)\to \Spec(S^{-1}A)
    \]
    is the restrition of $f^*$ to $S^{-1}Y$, and $S^{-1}Y={f^*}^{-1}(S^{-1}X)$.
  \end{lemma}

  \begin{proof}
    We have the following commutative diagram.
    \[
      \begin{tikzcd}
        A \ar[r,"f"]\ar[d,"\phi_A"'] & B\ar[d,"\phi_B"] \\
        S^{-1}A\ar[r,"S^{-1}f"']& S^{-1}B
      \end{tikzcd}
    \] 
   Applying the contravariant functor $\Spec$, we get a diagram of topological spaces.
   \[
     \begin{tikzcd}
        X & Y\ar[l,"f^*"']\\
        \Spec(S^{-1}A)\ar[u,"\phi_A^*"] 
                      & \Spec(S^{-1}B)\ar[u,"\phi_B^*"']\ar[l,"S^{-1}f^{*}"]
     \end{tikzcd}
   \] 
   Then 
   \[
     f^*(S^{-1}Y) = \phi_A^*(S^{-1}f^*(\Spec(S^{-1}B))\subset S^{-1}X
   \]
   allows us to restrict to the following diagram.
   %\[
   %  \begin{tikzcd}
   %     S^{-1}X \ar[d,shift left = 1ex,"(\phi_A^*)^{-1}"] 
   %     & S^{-1}Y\ar[l,"f^*|_{S^{-1}Y}"']
   %     \ar[d, shift right = 1ex,"(\phi_B^*)^{-1}"'] \\
   %     \Spec(S^{-1}A)\ar[u,shift left = 1ex,"\phi_A^*"]
   %                   & \Spec(S^{-1}B)\ar[u,shift right = 1ex,"\phi_B^*"']\ar[l,"S^{-1}f^{*}"]
   %  \end{tikzcd}
   %\] 
   \[
     \begin{tikzcd}
        S^{-1}X & S^{-1}Y\ar[l,"f^*|_{S^{-1}Y}"'] \\
        \Spec(S^{-1}A)\ar[u,"\phi_A^*"]
                      & \Spec(S^{-1}B)\ar[u,"\phi_B^*"']\ar[l,"S^{-1}f^{*}"]
     \end{tikzcd}
   \] 
   By the preceding lemma, $\phi_A^*$ and $\phi_B^*$ are heomeomorphisms onto their image, whence the result follows.
  \end{proof}

  One last lemma precedes the proof of the desired statement.

  \begin{lemma}
    Let $\a$ be an ideal of $A$ and let $\b=\a B$ be its extension in $B$.
    Let $f:A/\a\to B/\b$ be the homomorphism induced by $f$.
    If $\Spec (A/\a)$ is identified with its canonical image $V(\a)\subset X$, and $\Spec (B/\b)$ with $V(\b)\subset Y$, show that $\ol{f}^*=f^*|_{V(\b)}$.
  \end{lemma}

  \begin{proof}
    We have the following commutative diagram.
    \[
      \begin{tikzcd}
        A \ar[r,"f"]\ar[d,"\pi_A"'] & B\ar[d,"\pi_B"] \\
        A/\a\ar[r,"\ol{f}"']& B/\b
      \end{tikzcd}
    \] 
    Once again applying Spec,
    \[
      \begin{tikzcd}
        X & Y\ar[l,"f^*"']\\
        \Spec(A/\a)\ar[u,"\pi_A^*"] 
          & \Spec(B/\b)\ar[u,"\pi_B^*"']\ar[l,"\ol{f}^*"]
      \end{tikzcd}
    \] 
    and restricting as before gives us our last diagram.
    \[
      \begin{tikzcd}
        V(\a) & V(\b)\ar[l,"f^*|_{V(\b)}"'] \\
        \Spec(A/\a)\ar[u,"\pi_A^*"]
              & \Spec(B/\b)\ar[u,"\pi_B^*"'] \ar[l,"\ol{f}^*"]
      \end{tikzcd}
    \] 
    By Exercise 2.18, $\pi_A^*$ and $\pi_B^*$ are homeomorphisms onto their images $V(\a)$ and $V(\b)$, so that the result follows.
  \end{proof}

  Finally, we set $X=\Spec B$, $Y= \Spec A$, and fix $y=\p\in \Spec A$.
  Clearly, we have homeomorphisms
  \begin{align*}
    X_\p&=\Spec (B\otimes_A k(\p)) \\
       &\cong \Spec(B\otimes_A A_\p/\m_\p) \\
       &\cong \Spec(B_{A_\p}/\p B_{A_\p}) \\
       &\cong V(\p B_{A_\p}),
  \end{align*}
  where the last equality follows from Exercise 2.18 as before.
  Then equalities
  \begin{align*}
       &\cong V(\vphi_\p (\m_\p)B_{A_\p}) \\
       &\cong (\vphi_\p^*)^{-1}(V(\m_\p)) \\
       &\cong f_\p^{-1}(\m_\p) \\
       &\cong f^{-1}(\p).
  \end{align*}
  follow from minimality of the extension of an ideal, maximality of $\m_\p\lideal A_\p$, and the second lemma respectively.\footnote{I originally wrote this out a while ago, so I hadn't written it up again, let alone in simplied form, when we had our meeting. For some reason I thought (b) was the interesting part lol.}

  %(a) The affine case is Ex. 3.20 in Atiyah-MacDonald. %\footnote{I did it on paper, and it's a tad boring to write out. Just a series of commutative diagrams, and a repeat argument from Ex 2.1 of Hartshorne.} 
  Fix affine open covers $\{U_i=\Spec A_i\}$ of $X$ and $\{V_j=\Spec B_j\}$ of $Y$.
  Let $g:\Spec k(y)\to Y$ be the natural map, giving $\Spec k(y)$ the structure of a $Y$-scheme.
  We have 
  \[
    k(y)= \O_{Y,y}/m_y\O_{Y,y}\cong (B_j)_y/y(B_j)_y := k_j(y)
  \] 
  for all $y\in Y_j$.

  Then $X_y$ is the gluing of the spaces $f^{-1}(Y_j)\times_{Y_j} \Spec k_j(y)$, over all $Y_j\ni y$.
  We can write 
  \begin{align*}
    X_i\cap f^{-1}(Y_j)= \bigcup_{\substack{g_{i j k}\in X_i \\ D_i(g_{i j k})\subset f^{-1}(Y_j)}} D_i(g_{ij k})
  \end{align*}
  as a union of affines.
  Defining $f_{i j k}:=f|_{D_i(g_{i j k})}:D_i(g_{ i j k})\to Y_j$, the affine case gives us that
  \begin{align*}
    X_y &= \bigcup_{\substack{i,\,j,\, g_{ij k}\in X_i \\ D_i(g_{i j k})\subset f^{-1}(Y_j)}}
    D_i(g_{i j k})\times_{Y_j} \Spec k_j(y) \\
        &= \bigcup_{i,\, j,\, k} D_i(g_{i j k})_y \\
        &= \bigcup_{i,\, j,\, k} f_{i j k}^{-1}(y) \\
        &= f^{-1}(y)\cap \bigcup_{i, \, j, \,k} D_i(g_{ij k}) \\
        &= f^{-1}(y).
  \end{align*}

  (b) We have that $X_a=\Spec A_a$, where
  \begin{align*}
    A_a &= \frac{k[s,t]}{(s-t^2)}\otimes_{k[s]}\frac{k[s]_{(s-a)}}{(s-a)k[s]_{(s-a)}} \\
        &\cong k[t]\otimes_{k[t^2]}\frac{k[t^2]}{(a-t^2)} \\
        &\cong \frac{k[t]}{(a-t^2)}.
  \end{align*}
  Then $\Spec A_a\cong V(a-t^2)=\{(\sqrt{a}+ t,\sqrt{a}-t)\}\subset \Spec k[t]$.
  If $a\neq 0$, then this is a two point set, with residue fields given by
  \begin{align*}
    k_a(\sqrt{a}-t) &=\frac{k[t]}{(\sqrt{a}-t)}\cong k,\\
    k_a(\sqrt{a}+t) &=\frac{k[t]}{(\sqrt{a}+t)}\cong k.
  \end{align*}
  If $a=0$, then $\Spec A_0=\{t\}$ is singleton, and $t\in A_0$ is a nilpotent.

  Finally, we consider 
  \begin{align*}
    A_\eta &\cong \frac{k[s,t]}{(s-t^2)}\otimes_{k[s]}\frac{k[s]_{(0)}}{(0)k[s]_{(0)}} \\
           &\cong k[t]\otimes_{k[t^2]}k(t^2) \\
           &\cong k(s)[\sqrt{s}].
  \end{align*}
  Then $A_\eta$ is a PID, so its only ideals are of the form $(f(\sqrt{s}))$ for $f(\sqrt{s})= f_1(s)+f_2(s)\sqrt{s}$.
  Then $f\ol{f}(\sqrt{s})\in k(s)$ defined in the obvious manner is invertible whenever $f\neq 0$, so $(f(\sqrt{s}))=0,A_\eta$ gives us that $\Spec A_\eta$ is singleton.
  Furthermore, $(A_\eta)_{(0)}\cong k(\sqrt{s})$, which is a degree two extension over $k(\eta)=k(s)$, as desired.
\end{proof}

% \begin{enumerate}
% 	\item[3.11.] {
%      $(*)$ \emph{Closed Subschemes.}
%      \begin{enumerate}
%        \item Closed immersions are stable under base extension: if $f:Y\to X$ is a closed immersion, and if $X'\to X$ is any morphism, then $f':Y\times_X X' \to X'$ is also a closed immersion.
%        \item If $Y$ is a closed subscheme of an affine scheme $X=\Spec A$, then $Y$ is also affine, and in fact $Y$ is the closed subscheme determined by a suitable ideal $\a\subset A$ as the image of the closed immersion $\Spec A/\a\to \Spec A$.
%          [\emph{Hints}: First show that $Y$ can be covered by a finite number of open affine subsets of the form $D(f_i)\cap Y$, with $f_i\in A$.
%          By adding some more $f_i$ with $D(f_i)\cap Y$ if necessary, show that we may assume the $D(f_i)$ cover $X$.
%          Next, show that $f_1,\cdots, f_r$ generate the unit ideal of $A$.
%          Then use Ex. 2.17b to show that $Y$ is affinem and Ex. 2.18d to show that $Y$ comes from an ideal $\a\subset A$]
%          \emph{Note}: We will prove this with sheaves of ideals later.
%        \item Let $Y$ be a closed subset of a scheme $X$, and give $Y$ the reduced induced sub-scheme structure.
%          If $Y'$ is any other closed subscheme of $X$ with the same underlying topological space, show that the closed immersion $Y\to X$ factors through $Y'$.
%          We express this property by saying that the reduced induced structure is the smallest subscheme structure on a closed subset.
%        \item Let $f:Z\to X$ be a morphism.
%          Then there is a unique closed subscheme $Y$ of $X$ with the following property:
%          the morphism $f$ factors through $Y$, and if $Y'$ is any other closed subscheme of $X$ through with $f$ factors, then $Y\to X$ factors throuogh $Y'$ also.
%          We call $Y$ the \emph{scheme-theoretic image} of $f$.
%          If $Z$ is a reduced scheme, then $Y$ is just the reduced induced structure on the closure of the image $f(Z)$.
%      \end{enumerate}
%     }
% \end{enumerate}
 
\begin{enumerate}
 	\item[3.12.] {
      \emph{Closed subschemes of $\Proj S$.}
      \begin{enumerate}
        \item Let $\vphi:S\to T$ be a surjective homomorphism of graded rings, preserving degrees.
          Show that the open set $U$ is equal to $\Proj T$, and the morphism $f:\Proj T\to \Proj S$ is a closed immersion.
        \item If $I\subset S$ is a homogeneous ideal, take $T=S/I$ and let $Y$ be the closed subscheme of $X=\Proj S$ defined as image of the closed immersion $\Proj S/I\to X$.
          Show that different homogeneous ideals can give rise to the same closed subscheme.
          For example, let $d_0$ be an integer and let $I'=\oplus_{d\geqs 0}I_d$.
          Show that $I$ and $I'$ determine the same subscheme.\\
          We will see later that every closed subscheme of $X$ comes from a homogeneous ideal $I$ of $S$ (at least in the case where $S$ is a polynomial ring over $S_0$).
      \end{enumerate}
    }
\end{enumerate}
 
\begin{proof}
  (a) Denote $Y=\Proj S$ and $X=\Proj T$.
  We have $\vphi(S_+)= T_+$, so that $U=X$ by definition.
  Fix homogeneous $g\in S_+$.
  Surjectivity of $\vphi:S\to T$ implies $\vphi_{g}:S_{g}\to T_{\vphi(g)}$ is surjective.
  Then $\vphi$ degree preserving gives us that $(S_g)_d=\vphi_g^{-1}((T_g)_d)$ for all $d\geqs 0$, whence $\vphi_g:S_{(g)}\to T_{(\vphi(g))}$ is surjective.
  It follows that $f^\sharp:\O_X\to f_*\O_Y$ is surjective.

  Surjectivity of $\vphi$ gives us $\vphi(\vphi^{-1}(\b))$ for all $\b\trianglelefteq T$,  whence
  \[
    \p\supset \b \iff \vphi^{-1}(\p)\supset \vphi^{-1}(\b).
  \] 
  Then
  \begin{align*}
    f(V(\b)) &= V(\vphi^{-1}(\b))\cap f(X)
  \end{align*}
  shows that, in particular,
  \[
    f(X) = V(\ker\vphi)\cap f(X)\subset V(\ker\vphi).
  \] 
  Take $\p\in V(\ker\vphi)$, and $a\in \vphi^{-1}(\vphi(\p))$.
  Then we have $a'\in \p$ with $\vphi(a)=\vphi(a')$, whence $a-a'\in \ker\vphi\subset \p$ implies $a\in \p$.
  Thus, $\p=\vphi^{-1}(\vphi(\p))\in f(X)$ gives us $f(X)=V(\ker\vphi)$.
  We conclude that 
  \[
    f(V(\b)) = V(\ker\vphi + \vphi^{-1}(\b))
  \] 
  is closed in $\Proj S$, so that $f$ is a homeomorphism onto the closed set $f(X)=V(\ker\vphi)$.

  (b) Let $\pi: S\to S/I$ and $\pi':S\to S/I'$ be the natural projections.
  Since $I'\subset I=\ker\pi$, we can factor $\pi=\vphi\pi'$, where $\vphi:S/I'\to S/I$ is the identity $\vphi_d:S_d/I'_d=S_d/I_d\to S_d/I_d$ for $d\geqs d_0$.
  Then Ex. 2.14(b) gives us that $\vphi^*:\Proj (S/I)\to \Proj (S/I')$ is an isomorphism of sheaves, so $(\pi')^*(\Proj (S/I'))=\vphi^*\pi^*(\Proj (S/I))\cong \pi^*(\Proj(S/I))$ give the same closed subscheme structure.
\end{proof}

\begin{enumerate}
 	\item[3.13.] {
       \emph{Properties of Morphisms of Finite Type}
       \begin{enumerate}
         \item A closed immersion is a morphism of finite type.
         \item A quasi-compact open immersion is of finite type.
         \item A composition of two morphisms of finite type is of finite type.
         \item Morphisms of finite type are stable under base extension.
         \item If $X$ and $Y$ are schemes of finite type over $S$, then $X\times_S Y$ is of finite type over $S$.
         \item If $X\xrightarrow{f}Y\xrightarrow{g}Z$ are two morphisms, and if $f$ is quasi-compact, and $g\circ f$ is of finite type, then $f$ is of finite type.
         \item If $f:X\to Y$ is a morphism of finite type, and if $Y$ is noetherian, then $X$ is noetherian.
       \end{enumerate}
     }
\end{enumerate}

\begin{proof}
  (a) Let $f:X\to Y$ be a closed immersion, $\{U_i=\Spec A_i\}_{i\in I}$ an affine open cover of $X$, and fix $V=\Spec B\subset Y$ open affine.
  Then
  \begin{align*}
    f^{-1}(V)&= \bigcup_{i\in I} f^{-1}(V)\cap U_i \\
             &= \bigcup_{\substack{i\in I,~g_{ij}\in A_i \\ D_i(g_{ij}) \subset f^{-1}(V)}} D_i(g_{ij})
  \end{align*}
  is a union of open affines $D_i(g_{ij})=\Spec (A_i)_{g_{ij}}$.

  We have that $f^\sharp:\O_Y\to f_*\O_X$ is surjective.
  That is, $f^\sharp_\q:B_\q\to A_{B_\q}$ is surjective for each prime ideal $\p$ of $B$.
  Then
  \begin{align*}
    f(D_{i}(g_{ij})) &= \{f^\sharp(V)^{-1}(\p)\in \Spec B:\p\lideal A_i, ~\p\not\ni g_{ij}\} \\
                       &= \{\}
  \end{align*}

  (d) Let $f:X\to Y$ be a scheme of finite type, and $Y\to Y'$ a map of schemes.
\end{proof}

% \begin{enumerate}
%   \item[3.14.] {
%       $(*)$ If $X$ is a scheme of finite type over a field, show that the closed points of $X$ are dense.
%       Give an example to show that this is note true for arbitrary schemes.
%     }
% \end{enumerate}
% 
% \begin{enumerate}
%   \item[3.15.] {
%       $(*)$ Let $X$ be a scheme of finite type over a field $k$ (not necessarily algebraically closed).
%       \begin{enumerate}
%         \item Show that the following three conditions are equivalent (in which case we say that $X$ is \emph{geometrically irreducible}).
%           \begin{enumerate}
%             \item $X\times_k \ol{k}$ is irreducible, where $\ol{k}$ is the algebraic closure of $k$.
%             \item $X\times_k k_s$ is irreducible, where $k_s$ denotes the sperable closure of $k$.
%             \item $X\times_k K$ is irreducible for every extensions $K/k$.
%           \end{enumerate}
%         \item Show that the following three conditions are equivalent (in which case we say $X$ is \emph{geometrically reduced}).
%           \begin{enumerate}
%             \item $X\times_k \ol{k}$ is reduced.
%             \item $X\times_k k_p$ is reduced, where $k_p$ denotes the perfect closure of $k$.
%             \item $X\times_k K$ is reduced for all extensions $K/k$.
%           \end{enumerate}
%           \item We say that $X$ is \emph{geometrically integral} if $X\times_k \ol{k}$ is integral.
%             Note that, by (a) and (b), this implies $X=X\times_k k$ is integral.
%             Give examples of integral schemes which are neither geometrically irreducible nor geometrically reduced.
%       \end{enumerate}
%     }
% \end{enumerate}
% 
% \begin{enumerate}
%   \item[3.20.] {
%       $(*)$ \emph{Dimension}. Let $X$ be an integral scheme of finite type over a field $k$ (not necessarily closed).
%       Use results about varieties to prove the following.
%       \begin{enumerate}
%         \item For any closed point $p\in X$, $\dim X=\dim \O_P$, where for rings we always mean the Krull dimension.
%         \item Let $K(X)$ be the function field of $X$.
%           Then $\dim X=\trd K(X)/k$.
%         \item If $Y$ is a closed subset of $X$, then $\codim(Y,X)=\inf{\dim \O_{X,p}|p\in Y}$.
%         \item If $Y$ is a closed subset of $X$, then $\dim Y+\codim (Y,X)=\dim X$.
%         \item If $U$ is a nonempty open subset of $X$, then $\dim U=\dim X$.
%       \end{enumerate}
%     }
% \end{enumerate}
% 
% \begin{proof}
% \end{proof}
% 
% \begin{enumerate}
%   \item[19.] {
%     }
% \end{enumerate}
% 
% \begin{proof}
% \end{proof}
% 
% \begin{enumerate}
%   \item[20.] {
%     }
% \end{enumerate}
% 
% \begin{proof}
% \end{proof}
% 
% \begin{enumerate}
%   \item[21.] {
%     }
% \end{enumerate}
% 
% \begin{proof}
% \end{proof}

\end{document}
